\section{Summary and Looking Ahead}
\label{sec:ch4-summary}

An ordinary differential equation converts a local law of change into a trajectory. First-order equations describe relaxation, transport between compartments, growth, and circuit response. Second-order equations describe inertia, oscillation, damping, and resonance. Systems of equations expose eigenmodes and coupled motion, while phase space replaces a single graph by a geometric flow of possible states. Equilibria, linearization, Lyapunov functions, and bifurcations distinguish persistence from instability and reveal how qualitative behavior changes with parameters.

Analytical formulas remain indispensable, but numerical integration is part of the mathematical theory rather than an afterthought. Step size, consistency, stability, stiffness, invariants, and error control determine whether a computed trajectory represents the differential equation or merely the algorithm.

The central pattern is
\[
\boxed{\begin{gathered}
\text{physical assumption}\longrightarrow\text{differential equation}
\longrightarrow\text{initial data}\\
\longrightarrow\text{evolution}\longrightarrow\text{interpretation}
\end{gathered}}
\]
This pattern will recur throughout the book. Newton's laws become dynamical systems; Maxwell's equations propagate fields; the Schr\"odinger equation evolves quantum states; and relativistic field equations determine the geometry and matter content of spacetime. The next chapter extends the mathematical foundation needed to analyze these theories with greater structural precision.
